%% kapa: exam=1 round=35 year=2024 period=1 form=A subject=real_estate subject_ko=부동산학원론 seq=7 number=87
%% source: Q-Net 감정평가사 기출문제 (cst00309) artlSeq=5213229 file=raw/1008/5213229/1교시 문제지 원본.pdf
%% conversion: kordoc 4.24.51 -> tools/gen1.py (auto); review=proofread-sonnet; flags=none
\begin{question}{87}
부동산 중개계약에 관한 설명으로 ( )에 들어갈 것으로 옳은 것은?
\begin{framedtext}
( ㄱ ): 중개의뢰인이 특정한 개업공인중개사를 정하여 그 개업공인중개사에게 한정하여 해당 중개대상물을 중개하도록 하는 중개계약 \par
( ㄴ ): 중개의뢰인이 해당 중개대상물의 중개를 불특정 다수의 개업공인중개사에게 의뢰하고 먼저 거래를 성사시킨 개업공인중개사에게 보수를 지급하는 중개계약 \par
\end{framedtext}
\begin{choices}
  \item ㄱ: 일반중개계약, ㄴ: 전속중개계약
  \item ㄱ: 일반중개계약, ㄴ: 공동중개계약
  \item ㄱ: 전속중개계약, ㄴ: 공동중개계약
  \item ㄱ: 공동중개계약, ㄴ: 일반중개계약
  \item ㄱ: 전속중개계약, ㄴ: 일반중개계약
\end{choices}
\end{question}
\answer{5}
